\section{Related Work}
\label{sec:related}

We organize prior work into four groups and then state precisely what this paper adds.
Table~\ref{tab:mechanisms} compares mechanisms as described by their authors; it contains no
performance numbers, because we did not run those systems (Section~\ref{sec:setup} describes the
re-implemented baselines we did run).

\subsection{Fine-grained access control and query modification in databases}
Least privilege is a design principle for protection systems~\cite{saltzer1975protection}.
Relational databases realize it with discretionary grants on tables and
columns~\cite{griffiths1976authorization} and role-based models~\cite{sandhu1996rbac,ferraiolo2001nist},
with attribute-based variants~\cite{hu2014abac}. Row-level restriction has a long history:
Stonebraker's query modification rewrites a user query by conjoining the predicates of the views a
user may see~\cite{stonebraker1975query}; Rizvi et al.\ distinguish the \emph{Truman} model, which
silently rewrites queries to a user's view, from the \emph{non-Truman} model, which accepts a query only
if it is answerable from authorized views~\cite{rizvi2004extending}; Kabra et al.\ show how predicate
pushdown can leak information past such rewrites~\cite{kabra2006redundancy}. Commercial systems expose
these ideas as Oracle Virtual Private Database policies~\cite{oraclevpd} and PostgreSQL row-security
policies and \texttt{security\_barrier} views~\cite{pgrls,pgsecbarrier}. Surveys summarize the
area~\cite{bertino2005database}. Information disclosure through views has been formalized
\cite{miklau2004formal}, and statistical-database work shows that aggregate answers over authorized data
can still reveal individual values~\cite{denning1979tracker,adam1989security}. All of this machinery
enforces a policy written by an administrator for a user or role; none of it derives the policy from a
task description.

\subsection{Purpose-based access and privacy-by-design databases}
Hippocratic databases attach purposes to data and limit disclosure to what a stated purpose
requires~\cite{agrawal2002hippocratic,lefevre2004limiting}; purpose-based access control models purposes
hierarchically and checks access purposes against intended purposes~\cite{byun2008purpose}. Contextual
integrity frames privacy as appropriate information flow for a context~\cite{barth2006contextual}, and
data-minimization requirements appear in regulation~\cite{gdpr2016}. Purposes in these systems are
declared from a fixed vocabulary and bound to stored data in advance. Our task scope is closest in spirit
to a per-request purpose, but it is derived automatically from free-text tasks and expressed directly as a
relational scope (columns, row predicates), not as purpose labels. Two agent systems share this per-task view of
context. Conseca generates a security policy for each task from trusted context and enforces it deterministically on
agent actions~\cite{tsai2025conseca}; AirGapAgent minimizes the user data given to an agent according to contextual
integrity for the task at hand~\cite{bagdasarian2024airgapagent}. TSE applies the same idea to database reach and
enforces it inside the DBMS rather than in the agent runtime. Outside academia, short-lived, narrowly scoped
credentials are standard practice; OAuth 2.0 token exchange, for example, issues a down-scoped token for a specific
downstream call~\cite{rfc8693}. TSE can be read as token exchange for relational data, with the scope inferred from
the task. Industry guidance now promotes \emph{task-based authorization} for agents: permissions created per task in
a relationship-based authorization service and deleted or expired when the task ends~\cite{openfga2026task,stabile2026intent}.
These systems authorize agent actions on application objects and leave the mapping from intent to permissions to
the developer; TSE derives that mapping for relational data automatically and enforces it inside the DBMS.

\subsection{Least privilege and authorization for LLM agents}
Indirect prompt injection lets content an agent reads redirect its actions~\cite{greshake2023not}, and
benchmarks measure this for tool-using agents~\cite{debenedetti2024agentdojo,zhan2024injecagent,zhang2025asb,liu2024formalizing}.
Defenses that constrain authority include task-centric, hierarchical permission models for agent
delegates (MiniScope~\cite{zhu2025miniscope}), policy languages that restrict tool calls and their
arguments (Progent~\cite{shi2025progent}), dynamically generated task-scoped permissions for mobile/MCP
agents (AgentSentry~\cite{cai2025agentsentry}), short-lived intent certificates that narrow a tool
manifest (IGAC~\cite{zhu2026igac}), and capability leases composed from task intent and context source and
enforced at tool and OS level (IntentCap~\cite{zheng2026intentcap}). Further work argues for classic
security principles in agent design~\cite{zhang2025principles,li2025vision,sharma2026ac4a,surapani2026authz},
isolates or restructures execution~\cite{wu2025isolategpt,kim2025pfi}, separates control and data flow
(CaMeL~\cite{debenedetti2026camel}), applies information-flow control (FIDES~\cite{costa2025fides}), or
studies over-privileged tool selection~\cite{yang2026lowerpriv}; emulation frameworks expose risky tool use
before deployment~\cite{ruan2024toolemu}. A complementary line hardens the model itself against injected
instructions, by marking untrusted input (spotlighting~\cite{hines2024spotlighting}) or by training the model to
ignore instructions in the data channel (StruQ~\cite{chen2025struq}, SecAlign~\cite{chen2025secalign}). These
defenses lower the probability that an injection succeeds; TSE limits what a successful injection can reach, so the
two compose, and we do not compare against them. A data-centric survey catalogs privacy
risks of agents~\cite{lahjouji2026know}. These systems decide \emph{which tools or actions} an agent may
use, sometimes with argument constraints. For a database agent the relevant tool is usually a single
query interface; once it is allowed, the data it can reach is the full entitlement of the credential it
holds. Our work operates at that remaining layer: the rows and columns reachable through an allowed tool.

\subsection{Text-to-SQL under access control}
Text-to-SQL benchmarks~\cite{yu2018spider,li2023bird} and methods based on schema
linking~\cite{wang2020ratsql,lei2020reexamining,li2023resdsql} and LLM prompting~\cite{pourreza2023dinsql,gao2024dailsql}
assume unrestricted access. Recent work adds access control: Fei et al.\ synthesize roles and policies for
existing benchmarks and measure RBAC violations~\cite{fei2026rbac}; Klisura et al.\ evaluate whether models
refuse correctly under table/column policies~\cite{klisura2026role}; PCC-SQL enforces a role's column-use
policy deterministically during decoding~\cite{miyamoto2026pccsql}; SecureMCP places RBAC checks and SQL
interceptors in front of an LLM data-access server~\cite{kim2026securemcp}; MaskSQL abstracts schema and
values sent to a remote model~\cite{abedini2025masksql}. Prompt-to-SQL injection is a related application
risk~\cite{pedro2025prompt}. These systems enforce, or test compliance with, the \emph{user's} policy.
They do not narrow access below that policy for the task at hand.

\subsection{What is new relative to query modification and fine-grained access control}
\label{sec:whatsnew}
We do not claim views, row-level security, query modification, or role-based enforcement as new; we use
them unchanged. Relative to query modification~\cite{stonebraker1975query} and the Truman and non-Truman
models~\cite{rizvi2004extending}, the policy we enforce is not given: it is \emph{inferred per task} from
natural language and must be validated empirically for sufficiency, because a wrong policy breaks the
task (Section~\ref{sec:model}). Relative to Oracle VPD and PostgreSQL RLS~\cite{oraclevpd,pgrls}, our
contribution is a broker that compiles an inferred scope into ephemeral roles and \texttt{security\_barrier}
views owned by the user's role, so that the database, not the model, bounds the result by the user's
entitlement even when inference is wrong. Relative to Hippocratic and purpose-based
databases~\cite{agrawal2002hippocratic,byun2008purpose}, purposes are neither declared in advance nor
attached to stored data; the scope exists only for the lifetime of one task. Relative to agent tool
scoping~\cite{zhu2025miniscope,shi2025progent,cai2025agentsentry,zhu2026igac,zheng2026intentcap}, we
narrow the data reachable \emph{through} an allowed tool. Relative to access-controlled
text-to-SQL~\cite{fei2026rbac,klisura2026role,miyamoto2026pccsql,kim2026securemcp}, we enforce a
task-minimal subset of the user's policy, and we add a bounded, audited widening protocol for when the
inferred subset is insufficient. The new parts are therefore (i) per-task inference of a relational scope
from natural language, (ii) its compilation to ephemeral database objects whose ownership caps them at the
user's entitlement, (iii) the widening protocol, and (iv) the measured security/utility trade-off. To our
knowledge, no prior system combines (i) and (ii) (novelty search documented in the artifact\footnote{Code, preregistration, run logs and analysis scripts: \artifacturl.}).

\begin{table*}[t]
\centering
\caption{Mechanism-level comparison as described by each work (no performance numbers; we did not run these systems).
``Data reach'': what limits the rows/columns an agent can read through an allowed query tool.}
\label{tab:mechanisms}
\footnotesize
\begin{tabular}{p{2.6cm}p{2.5cm}p{2.7cm}p{2.6cm}p{2.4cm}p{2.6cm}}
\hline
\textbf{Work} & \textbf{Unit of authority} & \textbf{Source of policy} & \textbf{Enforcement point} & \textbf{Per-task?} & \textbf{Data reach via query tool} \\
\hline
Query modification~\cite{stonebraker1975query}, Truman/non-Truman~\cite{rizvi2004extending} & views / predicates & administrator & DBMS rewrite & no & user's views \\
Oracle VPD~\cite{oraclevpd}, PostgreSQL RLS~\cite{pgrls} & row predicates, column grants & administrator & DBMS & no & user's entitlement \\
Hippocratic DB~\cite{agrawal2002hippocratic}, PBAC~\cite{byun2008purpose} & purpose labels & declared purposes & DBMS & per declared purpose & data allowed for the purpose \\
MiniScope~\cite{zhu2025miniscope}, AgentSentry~\cite{cai2025agentsentry} & tools / permissions & task-centric roles; generated per task & agent runtime & yes & unchanged once tool allowed \\
Progent~\cite{shi2025progent}, IGAC~\cite{zhu2026igac} & tool calls (+arguments) & policy DSL; intent certificate & tool-call checker & yes (IGAC) & unchanged once query allowed \\
IntentCap~\cite{zheng2026intentcap} & capability leases & task intent + context sources & tool/OS IFC & yes & not relational \\
CaMeL~\cite{debenedetti2026camel}, FIDES~\cite{costa2025fides} & control/data flow, labels & trusted query; labels & interpreter / IFC & yes & unchanged \\
RBAC text-to-SQL~\cite{fei2026rbac,klisura2026role} & roles (tables/columns) & synthesized roles & evaluated (model compliance) & no & user's entitlement \\
PCC-SQL~\cite{miyamoto2026pccsql} & column-use policy & role policy & constrained decoding & no & role's columns \\
SecureMCP~\cite{kim2026securemcp} & RBAC + SQL checks & role policy & MCP server pipeline & no & role's entitlement \\
\textbf{TSE (this paper)} & rows, columns & \textbf{inferred from task} $\cap$ entitlement & \textbf{DBMS} (ephemeral role, views) & \textbf{yes} & \textbf{task scope (+bounded widening)} \\
\hline
\end{tabular}
\end{table*}
